\newpage
\newgeometry{top=1.5cm, bottom=1.5cm, left=1cm, right=1cm} % Margini ridotti per Cheat Sheet
\section*{Cheat Sheet Finale: Algebra Lineare}
\addcontentsline{toc}{section}{Cheat Sheet Finale}

\begin{multicols*}{2}

% --- SISTEMI LINEARI ---
\begin{tcolorbox}[colback=lightBg, colframe=cherryRed, title={Sistemi Lineari ($Ax=b$)}]
\small
\textbf{Rouché-Capelli:}
\begin{itemize}[leftmargin=*]
    \item Compatibile $\iff \text{rk}(A) = \text{rk}(A|b)$.
    \item $\infty^{n-r}$ soluzioni se $r < n$.
    \item 1 soluzione unica se $r = n$.
\end{itemize}
\textbf{Gauss:}
\begin{enumerate}[leftmargin=*]
    \item Scambio righe ($R_i \leftrightarrow R_j$).
    \item Moltiplicazione per scalare ($\lambda R_i$).
    \item Comb. lineare ($R_i \to R_i + k R_j$).
\end{enumerate}
\textbf{Cramer:} Se $\det(A) \neq 0, x_i = \frac{\det(A_i)}{\det(A)}$.
\end{tcolorbox}

% --- SPAZI VETTORIALI ---
\begin{tcolorbox}[colback=white, colframe=maroon, title={Spazi Vettoriali e Basi}]
\small
\textbf{Indipendenza Lineare:} $\sum \lambda_i v_i = 0 \implies \forall \lambda_i = 0$.
\textbf{Generatori:} $\text{Span}(v_1, \dots, v_k) = V$.
\textbf{Base:} Insieme di generatori lin. indip.
\textbf{Dimensione:} Numero elementi della base.
\tcblower
\textbf{Formula di Grassmann:}
\[ \dim(U+W) = \dim(U) + \dim(W) - \dim(U \cap W) \]
\textbf{Somma Diretta ($V = U \oplus W$):}
\begin{itemize}[leftmargin=*]
    \item $U \cap W = \{0\}$.
    \item Ogni $v \in V$ si scrive in modo unico come $u+w$.
\end{itemize}
\end{tcolorbox}

% --- MATRICI ---
\begin{tcolorbox}[colback=lightBg, colframe=cherryRed, title={Matrici e Determinante}]
\small
\textbf{Rango (rk):} Numero pivot / Numero righe lin. indip.
\textbf{Inversa ($A^{-1}$):} Esiste $\iff \det(A) \neq 0$.
\[ A^{-1} = \frac{1}{\det(A)} (\text{Cof}(A))^T \]
\textbf{Proprietà Determinante:}
\begin{itemize}[leftmargin=*]
    \item $\det(A) = \det(A^T)$.
    \item $\det(AB) = \det(A)\det(B)$ (Binet).
    \item $\det(A^{-1}) = 1/\det(A)$.
\end{itemize}
\end{tcolorbox}

% --- APPLICAZIONI LINEARI ---
\begin{tcolorbox}[colback=white, colframe=maroon, title={Applicazioni Lineari ($f: V \to W$)}]
\small
\textbf{Nucleo ($\ker f$):} $\{ v \in V \mid f(v) = 0 \}$. (Iniettiva $\iff \ker f = \{0\}$).
\textbf{Immagine ($\text{Im } f$):} $\{ w \in W \mid \exists v: f(v)=w \}$. (Suriettiva $\iff \text{Im } f = W$).
\tcblower
\textbf{Teorema della Dimensione (Rank-Nullity):}
\[ \dim(V) = \dim(\ker f) + \dim(\text{Im } f) \]
\textbf{Matrice Associata ($M_\mathcal{B}^\mathcal{C}(f)$):}
Le colonne sono le coordinate delle immagini dei vettori della base di partenza rispetto alla base di arrivo.
\[ [f(v)]_\mathcal{C} = A [v]_\mathcal{B} \]
\end{tcolorbox}

% --- AUTOVALORI ---
\begin{tcolorbox}[colback=lightBg, colframe=cherryRed, title={Autovalori e Diagonalizzazione}]
\small
\textbf{Definizione:} $f(v) = \lambda v \implies (A - \lambda I)v = 0$.
\textbf{Polinomio Caratteristico:} $P_A(\lambda) = \det(A - \lambda I)$.
\textbf{Spettro:} Insieme delle radici di $P_A(\lambda)$.
\textbf{Molteplicità:}
\begin{itemize}[leftmargin=*]
    \item Alg. ($m_a(\lambda)$): esponente nel pol. caratt.
    \item Geom. ($m_g(\lambda)$): $\dim(\ker(A-\lambda I)) = n - \text{rk}(A-\lambda I)$.
\end{itemize}
\textbf{Diagonalizzabile} $\iff$:
\begin{itemize}[leftmargin=*]
    \item Tutte le radici sono nel campo (p.c. si spezza).
    \item $\forall \lambda, m_g(\lambda) = m_a(\lambda)$.
\end{itemize}
Se diag., $D = P^{-1}AP$ con $P$ matrice cambio base (autovettori).
\end{tcolorbox}

% --- PRODOTTI SCALARI ---
\begin{tcolorbox}[colback=white, colframe=maroon, title={Prodotti Scalari e Spettale}]
\small
\textbf{Standard in $\mathbb{R}^n$:} $\langle x, y \rangle = x^T y$.
\textbf{Norma:} $\|v\| = \sqrt{\langle v, v \rangle}$.
\textbf{Ortogonalità:} $\langle v, w \rangle = 0$.
\textbf{Proiezione Ortogonale su $v$:} $P_v(u) = \frac{\langle u, v \rangle}{\|v\|^2} v$.
\textbf{Gram-Schmidt:} $u_{k} = v_k - \sum_{j=1}^{k-1} P_{u_j}(v_k)$.
\tcblower
\textbf{Teorema Spettrale (Reale):}
Una matrice $A \in \mathbb{R}^{n \times n}$ è simmetrica ($A=A^T$) $\iff$ ammette una base ORTONORMALE di autovettori.
(Tutte le matrici simmetriche sono diagonalizzabili).
\end{tcolorbox}

\end{multicols*}
\restoregeometry
